\chapter{Detection, Tracking, and Trajectory Generation}
\label{chap:detection}

\section{Overview}

This chapter describes how the system transforms raw drone video into structured metric
trajectories.  The output of this stage --- a CSV file containing one row per
(frame, track\_id) pair with real-world coordinates and velocity --- is the primary input
to all subsequent stages.

\section{Object Detection}

\subsection{RT-DETRv2-X with DINO Decoder (Detection Backbone)}

The detection backbone is RT-DETRv2-X, an Extra-Large variant of the Real-Time Detection
Transformer that retains the DINO (DETR with Improved deNoising Optimisation) decoder head.
The DINO decoder incorporates two mechanisms that are particularly beneficial for dense
Indian traffic:

\begin{itemize}
    \item \textbf{Contrastive denoising training (CDN):}
          During training, each object query learns to distinguish clean anchor assignments
          from noisy ones.  The result is tighter, more calibrated confidence scores
          per frame.  When two overlapping tiles detect the same motorcycle, the
          DINO-calibrated scores are close to each other (e.g.\ 0.60 and 0.59) rather than
          arbitrarily divergent (e.g.\ 0.82 and 0.43), which makes the downstream
          Weighted Box Fusion step far more reliable.

    \item \textbf{Multi-scale deformable attention:}
          The encoder explicitly attends to feature maps at four scales (P3--P6),
          enabling simultaneous detection of small motorcycles ($\approx 20$--50 px) and
          large buses ($\approx 200$+ px) within the same crowded frame.
\end{itemize}

Five COCO class IDs are retained: person~(0), car~(2), motorcycle~(3), bus~(5), and
truck~(7).  A global confidence threshold of $\text{conf} = 0.10$ is used at the SAHI
level; class-specific thresholds are applied in a separate post-processing step.

\subsection{Why DINO: Architectural Intuition and Choice Rationale}

Standard object detectors present clear trade-offs when applied to high-altitude aerial drone video:
\begin{itemize}
    \item \textbf{CNN Backbones (YOLOv8 / Faster R-CNN):} Traditional one-stage CNN detectors downsample feature maps by $32\times$, causing small two-wheelers ($< 32 \times 32$ pixels) to lose spatial activation. Furthermore, standard NMS produces uncalibrated confidence scores across overlapping image tiles, causing duplicate detections across SAHI sub-tiles to receive wildly divergent scores.
    \item \textbf{Standard DETR:} Vanilla DETR suffers from slow bipartite matching convergence during training and struggles with multi-scale feature alignment.
\end{itemize}

RT-DETRv2-X with a DINO decoder was chosen to address these specific failure modes:
\begin{enumerate}
    \item \textbf{Contrastive Denoising Training (CDN):} DINO explicitly trains decoder queries against contrastive noisy bounding-box anchors. This produces tightly calibrated confidence scores per frame, ensuring that duplicate detections across adjacent SAHI tiles yield consistent confidence scores (e.g., $0.60$ vs $0.59$), which enables smooth Weighted Box Fusion (WBF) deduplication.
    \item \textbf{Deformable Multi-Scale Attention:} The DINO encoder samples feature maps across four resolution scales (P3--P6) simultaneously, enabling high-recall detection of small two-wheelers ($20$--$50$\,px) and large buses ($>200$\,px) within dense, un-lane-disciplined traffic streams.
\end{enumerate}

\subsection{SAHI Tiling}

Standard inference on a full drone frame at the native resolution causes vehicles to
subtend very few pixels in the network's input tensor after downsampling.  SAHI (Slicing
Aided Hyper Inference) resolves this by dividing the frame into overlapping tiles before
inference, so that small vehicles appear at a much higher relative resolution.

The detector is wrapped in the \texttt{SahiDinoDetector} class
(\texttt{sahi\_dino\_detection.py}).  The tiling configuration used is:
\begin{itemize}
    \item Tile size: $384 \times 384$ pixels
    \item Overlap ratio: 0.75 (both dimensions)
    \item Post-SAHI fusion: Weighted Box Fusion (WBF) with cluster distance 45 px and
          IoU threshold 0.30
\end{itemize}

\subsection{Weighted Box Fusion}

After SAHI assembles predictions from all tiles, the \texttt{weighted\_box\_fusion()}
function (\texttt{sahi\_fusion.py}) is applied as a second deduplication pass beyond
SAHI's built-in NMS.  This eliminates low-IoU duplicate detections in the
$0.22$--$0.33$ IoU range that bypass standard NMS thresholds but would otherwise confuse
the ByteTrack Stage-1 matching.

\section{Object Tracking}

\subsection{ByteTrack}

The tracker (\texttt{tracker.py}) follows the ByteTrack framework~\cite{bytetrack}.
The per-frame update proceeds in two stages:

\begin{enumerate}
    \item \textbf{Stage 1:}
          High-confidence detections ($\text{conf} \geq 0.5$) are associated with all
          currently Tracked and recently Lost tracks using a cost matrix combining DIoU
          and appearance distance.  The Hungarian algorithm finds the minimum-cost
          assignment.

    \item \textbf{Stage 2:}
          Low-confidence detections are given a second attempt to associate with the tracks
          that remained unmatched after Stage~1.  This recovers tracks for partially occluded
          vehicles whose detection confidence has dropped temporarily.
\end{enumerate}

Each track maintains a constant-velocity Kalman filter with 8-dimensional state
$[c_x, c_y, w, h, \dot{c}_x, \dot{c}_y, \dot{w}, \dot{h}]^\top$, where $(c_x, c_y)$ is
the box centre and $(w, h)$ is its size.  Process noise is scaled proportionally to the
bounding-box area so that larger vehicles receive higher noise, reflecting their larger
absolute displacements.

\subsection{Appearance Feature Matching}

To reduce identity switches during occlusion, each bounding-box crop is processed by an appearance feature extractor (\texttt{reid.py}), which produces a 128-dimensional embedding using a ResNet50 backbone (classification head replaced by a linear projection).  A 512-bin 3-D HSV colour histogram is appended to improve discrimination of visually similar vehicles. Each track maintains an exponential moving-average appearance prototype:
\begin{equation}
    \hat{\mathbf{e}}_t = \alpha\,\hat{\mathbf{e}}_{t-1} + (1-\alpha)\,\mathbf{e}_t,
    \qquad \alpha = 0.9
    \label{eq:ema}
\end{equation}
Cosine distance between prototypes supplements the IoU cost in Stage-1 matching.

\subsection{Trajectory Smoothing}

The bounding-box centres exported per frame still carry quantisation noise and
detection jitter even after Kalman filtering.  The \texttt{MovingAverageSmoother}
(\texttt{smoothing.py}) applies a per-track simple moving average with window $W = 7$
frames:
\begin{equation}
    \tilde{c}(t) = \frac{1}{W}\sum_{k=0}^{W-1} c(t-k)
    \label{eq:sma}
\end{equation}
This introduces a maximum lag of 3 frames ($\approx 100$ ms at 30 FPS), which is negligible
relative to the safety-rule evaluation windows (90--1500 frames).

\section{World-Coordinate Calibration}
\label{sec:calibration}

\subsection{Motivation}

Pixel coordinates do not carry physical meaning.  A separation of 50 pixels between two
vehicle centres could correspond to 2 m or 20 m depending on the camera altitude, focal
length, and tilt.  Since all safety rules and ML features are defined in metres and metres
per second, a calibrated pixel-to-metre mapping is an essential prerequisite.

\subsection{Physical Measurement Procedure}

The calibration is grounded in a direct physical measurement of the roundabout.  The team
visited the actual intersection and measured the lane width perpendicular to the
circulatory direction using a laser measuring device.  The measured lane width is
$L_{\text{real}} = 7.0$ m.

The pixel extent of the same lane boundary was measured in the drone video frame using the
interactive calibration tool (\texttt{interactive\_calibration.py}), which allows the user
to click two points on the video frame and computes the pixel distance.  The measured
pixel span of the 7.0 m lane width is $L_{\text{px}} = 85$ pixels.

The uniform scale factor is therefore:
\begin{equation}
    s = \frac{L_{\text{real}}}{L_{\text{px}}} = \frac{7.0}{85} \approx 0.0824\;\text{m/px}
    \label{eq:scale}
\end{equation}

\subsection{Coordinate Transformation}

Each smoothed bounding-box pixel centre $(c_x, c_y)$ is converted to world coordinates
$(x_w, y_w)$ by applying the uniform scale factor:
\begin{equation}
    x_w = c_x \cdot s, \qquad y_w = c_y \cdot s
    \label{eq:world}
\end{equation}

The roundabout centre in world coordinates, computed by applying the same scale to the
pixel-space centre $(870, 570)$, is:
\begin{equation}
    (X_C, Y_C) = (870 \cdot s,\; 570 \cdot s) \approx (43.5,\; 28.5)\;\text{m}
    \label{eq:centre}
\end{equation}

The inner circulatory boundary ($120\,\text{px} \times s$) is $R_{\text{inner}} \approx 6.0$ m
and the outer boundary ($280\,\text{px} \times s$) is $R_{\text{outer}} \approx 14.0$ m.

Per-frame velocity is estimated as:
\begin{equation}
    v(t) = \sqrt{(x_w(t) - x_w(t{-}1))^2 + (y_w(t) - y_w(t{-}1))^2} \times f_{\text{ps}}
    \label{eq:velocity}
\end{equation}
where $f_{\text{ps}} = 30$ FPS.

\section{Output CSV Schema}

\begin{table}[H]
\centering
\caption{Tracking pipeline output CSV column schema.}
\label{tab:csv_schema}
\begin{tabular}{lll}
\toprule
\textbf{Column} & \textbf{Type} & \textbf{Description} \\
\midrule
\texttt{frame}              & int   & 0-based video frame index \\
\texttt{track\_id}          & int   & Persistent unique identifier \\
\texttt{class\_name}        & str   & Road-user class (car, motorcycle, \ldots) \\
\texttt{x1, y1, x2, y2}    & float & Bounding box in pixel coordinates \\
\texttt{center\_x, center\_y} & float & Smoothed pixel centre \\
\texttt{world\_x, world\_y}   & float & World coordinates (metres) \\
\texttt{confidence}         & float & Detector score $\in [0,1]$ \\
\texttt{velocity\_ms}       & float & Estimated ground speed (m/s) \\
\bottomrule
\end{tabular}
\end{table}
